% !TeX program = xelatex
% 教辅《三维设计》4.3.1，第62—63页“师生共研”。按纸本题号编排。
% 40分钟：唤起定义3分钟，典例1及对点练习10分钟，典例2及辨析10分钟，
% 典例3共13分钟，出口题与回看4分钟。出口题之后为选讲。
% 每帧首次只投任务。巡视或听完口答后用右箭头揭示；左箭头回到上一步。
\documentclass[aspectratio=169,11pt]{beamer}
\usepackage[UTF8,fontset=fandol]{ctex}
\IfFontExistsTF{Microsoft YaHei}{\setCJKsansfont{Microsoft YaHei}}{}
\usetheme{GaoZhongYuan}
\hypersetup{pdfpagemode=FullScreen}
\setbeamercovered{invisible}
\setlength{\parskip}{.55em}
\newcommand{\wherein}[1]{{\small\color{campusgreen}\bfseries #1}\par\medskip}
\newcommand{\ask}[1]{\par\medskip{\bfseries #1}\par}
\newcommand{\work}[1]{\par\smallskip{\color{campusgreen}#1}\par}
\title{4.3.1 对数的概念}
\author{深圳中学高中园}
\date{教辅第62—63页}
\begin{document}
\campusmaketitle

\begin{frame}[t]{对数记下的是哪个数？}
\wherein{第62页 · 翻到“师生共研”}
\[ 2^3=8 \qquad\qquad 2^x=10 \]
\ask{左式中的指数是几？右式的指数介于哪两个整数之间？}
\par\pause\par
$2^3<10<2^4$，所以 $3<x<4$。把这个指数记作 $\log_2 10$。
\par\pause\par
\work{当 $a>0$、$a\ne1$、$N>0$ 时，\quad
$a^x=N\ \Longleftrightarrow\ x=\log_a N$。}
指出两种写法中的底数；幂 $N$ 在对数式中叫什么？
\end{frame}

\begin{frame}[t]{哪些取值能使对数有意义？}
\wherein{第62页 · 典例1（1）}
在 $b=\log_{a-2}(5-a)$ 中，实数 $a$ 的取值范围是（\quad）。
\[
\begin{array}{ll}
\mathrm A.\;(-\infty,2)\cup(5,+\infty) & \mathrm B.\;(2,5)\\[2pt]
\mathrm C.\;(2,3)\cup(3,5) & \mathrm D.\;(3,4)
\end{array}
\]
\ask{圈出底数和真数，在书上写全限制条件，再举手报出选项。}
\par\pause\par
底数给出 $a-2>0,\ a-2\ne1$；真数给出 $5-a>0$。
\par\pause\par
\work{取交集得 $(2,3)\cup(3,5)$，选 C。若选 B，代入 $a=3$ 检查底数。}
\end{frame}

\begin{frame}[t]{底数、指数和幂各去了哪里？}
\wherein{第62页 · 典例1（2）}
\[
\begin{array}{ll}
(1)\;5^3=125 & (2)\;\log_2 16=4\\[5pt]
(3)\;10^{-2}=0.01 & (4)\;\log_{\sqrt5}125=6
\end{array}
\]
\ask{在书上完成四个互化式。同桌互读，检查每一式的底数。}
\par\pause\par
(1) $\log_5 125=3$ \hfill (2) $2^4=16$\hspace{3cm}\par
\par\pause\par
(3) $\lg 0.01=-2$ \hfill (4) $(\sqrt5)^6=125$\hspace{2cm}\par
\work{$\lg$ 的底数是 $10$。第(3)题中，负数表示哪个量？}
\end{frame}

\begin{frame}[t]{两个等式都成立，就算互化正确吗？}
\wherein{第63页 · 题型一 · 对点练清 2（多选）}
选出互化正确的各组，并把其余组改正确。
\[
\begin{array}{lll}
\mathrm A.&10^0=1 & \lg 1=0\\
\mathrm B.&27^{-1/3}=1/3 & \log_{27}(1/3)=-1/3\\
\mathrm C.&\log_3 9=2 & 9^{1/2}=3\\
\mathrm D.&\log_5 5=1 & 5^1=5
\end{array}
\]
\par\pause\par
\work{选 A、B、D。C 的两个等式各自成立；$\log_3 9=2$ 直接化成 $3^2=9$。}
\end{frame}

\begin{frame}[t]{先看未知数在什么位置}
\wherein{第63页 · 典例2（1）（2）}
\[ (1)\ \lg 0.01=x \qquad\qquad (2)\ \log_7(x+2)=2 \]
\ask{各写一个指数式。请两位同学板演第一行，并说明自己要求哪个量。}
\par\pause\par
$(1)\quad 10^x=0.01=10^{-2}$，所以 $x=-2$。
\par\pause\par
$(2)\quad 7^2=x+2$，所以 $x=47$。\par
\work{把 $47$ 代回第（2）题，真数是 $49$，恰好为 $7^2$。}
\end{frame}

\begin{frame}[t]{底数小于 1，怎样判断指数的符号？}
\wherein{第63页 · 典例2（3）（4）}
\[ (3)\ \log_{2/3}\frac94=x \qquad\qquad (4)\ x=\log_{1/2}32 \]
\ask{先预测两个答案的符号，再把真数写成相应底数的幂。}
\par\pause\par
$(3)\quad \dfrac94=\left(\dfrac32\right)^2
=\left(\dfrac23\right)^{-2}$，故 $x=-2$。
\par\pause\par
$(4)\quad 32=2^5=\left(\dfrac12\right)^{-5}$，故 $x=-5$。\par
\work{两题的底数均在 $(0,1)$ 内，真数都大于 $1$，所以指数为负。}
\end{frame}

\begin{frame}[t]{同样得到平方方程，负根都能保留吗？}
\wherein{第63页 · 题型二 · 对点练清（2）（3）}
\[ (2)\ \log_x 25=2 \qquad\qquad (3)\ \log_5(x^2)=2 \]
\ask{先独立求解，再与同桌解释每个负根的取舍。}
\par\pause\par
两题都推出 $x^2=25$。代入 $x=-5$，各自的底数和真数是多少？
\par\pause\par
\work{（2）底数 $x>0$ 且 $x\ne1$，只取 $5$。\\
（3）真数是 $x^2$，取 $5$ 或 $-5$ 都有 $x^2=25>0$。}
\end{frame}

\begin{frame}[t]{三道计算，先做哪一步？}
\wherein{第63页 · 典例3（1）}
\[ (1)\ 2^{-\log_2 3}\qquad (2)\ e^{3\ln 7}\qquad (3)\ \lg(0.001^2) \]
\ask{分别标出负号、系数 3 和平方的位置。同桌各选一道讲第一步。}
\par\pause\par
由定义，$a^{\log_a N}=N$（$a>0$，$a\ne1$，$N>0$）。\par (1) $2^{-\log_2 3}=\dfrac{1}{2^{\log_2 3}}=\dfrac13$。
\par\pause\par
(2) $e^{3\ln 7}=(e^{\ln 7})^3=7^3=343$，其中 $\ln 7=\log_e 7$。
\par\pause\par
(3) $\lg(0.001^2)=\lg((10^{-3})^2)=\lg(10^{-6})=-6$。
\end{frame}

\begin{frame}[t]{把哪一整块当作真数？}
\wherein{第63页 · 典例3（2）(1)}
\[ \log_3(\lg x)=1 \]
\ask{圈出最外层的真数。只写第一步，然后停笔。}
\par\pause\par
$\lg x=3^1=3$。\qquad 下一步要把哪个数作底？
\par\pause\par
\work{$x=10^3=1000$。验算时从里往外算，$\lg 1000=3$，$\log_3 3=1$。}
\end{frame}

\begin{frame}[t]{三层对数，一次解开一层}
\wherein{第63页 · 典例3（2）(2)}
\[ \log_3\!\bigl[\log_4(\log_5 x)\bigr]=0 \]
\ask{请一位同学板演，其他同学在书上写。每一步给下一位同学留一个等式。}
\par\pause\par
$\log_4(\log_5 x)=3^0=1$
\par\pause\par
\par $\log_5 x=4^1=4$
\par\pause\par
\par $x=5^4=625$。\qquad 从里向外验算，依次得到 $4,\ 1,\ 0$。
\end{frame}

\begin{frame}[t]{下课前，独立写下两行}
\wherein{第62页 · 题型一 · 对点练清 1}
使 $\log_a(1-2a)$ 有意义的 $a$ 的取值范围是什么？
\ask{第一行列全条件，第二行写交集。和典例1相比，这次还要挖去一个点吗？}
\par\pause\par
$a>0,\quad a\ne1,\quad 1-2a>0$。
\par\pause\par
\work{$0<a<\tfrac12$，选 B。这个区间已经排除了 $1$。}
\medskip
合上书口述：$a^x=N$ 中的指数怎样用对数表示？
\end{frame}

\begin{frame}[t]{选讲：把右端的 0 改成 1}
\wherein{第63页 · 题型三 · 对点练清 1}
\[ \log_3\!\bigl[\log_4(\log_5 x)\bigr]=1 \]
\ask{仿照典例3写出三步。最后一行用幂表示即可。}
\par\pause\par
$\log_4(\log_5 x)=3^1=3$，进而 $\log_5 x=4^3=64$。
\par\pause\par
\work{$x=5^{64}$。从里往外算，三个对数的值依次为 $64,\ 3,\ 1$。}
\end{frame}

\begin{frame}[t]{选练：指数与真数}
\wherein{第63页 · 题型二 · 对点练清（1）（4）}
\[ (1)\ \log_2 x=-\tfrac12 \qquad (4)\ 2\log_3 x=4 \]
\ask{各写一个能直接读出真数的指数式。}
\par\pause\par
$(1)\quad x=2^{-1/2}=\dfrac{\sqrt2}{2}$。
\par\pause\par
$(4)\quad \log_3 x=2$，故 $x=3^2=9$。
\end{frame}

\begin{frame}[t]{选练：底数换了，恒等式还怎么用？}
\wherein{第63页 · 题型三 · 对点练清 2、3}
（2）沿用典例3（2）的 $x=625$，求 $625^{\log_x 3}$。\\[5pt]
（3）若 $3^{\log_3[\log_4(\log_5 x)]}=1$，求 $x$。
\ask{先圈出能配对的底数，再分别化简。}
\par\pause\par
（2）$625^{\log_{625}3}=3$。
\par\pause\par
（3）$\log_4(\log_5 x)=1$，故 $\log_5 x=4$，$x=625$。\\
代回时外层对数的真数为 $1$，满足定义条件。
\end{frame}
\end{document}
